\subsection{Conceptos escolares básicos}
\label{subsec:conceptos_escolares}

Para entender con claridad el objetivo de la página web, es importante conocer cómo se vive el proceso de tesis en la universidad y cuáles son las situaciones cotidianas que pueden retrasar a los estudiantes.

\subsubsection{El proceso de tesis en el posgrado}
Realizar una tesis de posgrado es un trabajo formativo formal donde el estudiante elige un problema científico, investiga soluciones posibles y redacta un reporte escrito con rigor y orden metodológico \parencite{sampieri2014}. A lo largo de este camino, el texto pasa por constantes correcciones, borradores y revisiones conjuntas con el profesor asesor antes de que el comité autorice la versión definitiva para el examen de grado.

\subsubsection{La comunicación entre el asesor y el alumno}
El acompañamiento tutorial frecuente entre el docente y el tesista es el factor más determinante para que el trabajo continúe con buen ritmo \parencite{sim2023technology}. Cuando las sesiones de trabajo se espacian demasiado o los acuerdos se toman únicamente en charlas informales sin dejarlos por escrito, las tareas se olvidan con facilidad y el estudiante suele perder el hilo de su proyecto.

\subsubsection{El rezago escolar en el posgrado}
El rezago escolar sucede cuando un estudiante sobrepasa los tiempos permitidos por el plan de estudios sin entregar su proyecto concluido o abandona la carrera por falta de supervisión continua \parencite{lando2025digital}. Contar con una plataforma oficial que revise de forma constante las fechas permite notar a tiempo este riesgo y brindar ayuda al alumno antes de que se quede sin opciones de titulación.